\section{Systematic Survey of 25 Studies (2023--2025)}
\label{sec:survey}

\subsection{Survey protocol}
The survey was carried out in the initial stage of the project and recorded in a shared workbook with the same fields for every paper: title, publication, abstract, techniques, results, year, link, research gaps and key takeaways. The work was divided among the three team members. The search was limited to peer-reviewed journal and conference papers from 2023 to 2025 on agricultural price prediction, volatility and anomaly analysis, hybrid and ensemble forecasting, and smart-farming decision support, drawn from Springer Nature, Elsevier ScienceDirect and IEEE Xplore. After removing duplicates (six IEEE papers were recorded by two members), \textbf{25 distinct studies} remained, labelled S1 to S25. Every study was matched to a bibliographic record in Crossref before it was cited. The findings below are the authors' own, as recorded in the workbook from the abstracts and results sections; performance figures were not re-measured by us.

\begin{table}[!t]
\centering
\caption{Distribution of the 25 surveyed studies by publisher and year (as recorded in the workbook)}
\label{tab:survey-dist}
\small
\begin{tabular}{|l|c|c|c|c|}
\hline
\textbf{Source} & \textbf{2023} & \textbf{2024} & \textbf{2025} & \textbf{Total} \\ \hline
Elsevier (ScienceDirect), S1--S10 & 2 & 5 & 3 & 10 \\ \hline
Springer Nature, S11--S19 & 2 & 5 & 2 & 9 \\ \hline
IEEE, S20--S25 & 1 & 3 & 2 & 6 \\ \hline
\textbf{Total} & \textbf{5} & \textbf{13} & \textbf{7} & \textbf{25} \\ \hline
\end{tabular}
\end{table}

\subsection{Overview of the studies}
Table~\ref{tab:survey-overview} lists the studies with their venue and the task they address. Of the 25, 15 forecast prices (S2 also detects anomalies), 5 concern volatility or its drivers, 1 is about anomaly detection, and 4 are crop-recommendation systems or reviews.

\Needspace{12\baselineskip}
{\footnotesize\setlength{\tabcolsep}{4pt}\renewcommand{\arraystretch}{1.15}
\begin{longtable}{|p{0.9cm}|p{7.7cm}|p{4.2cm}|p{2.7cm}|}
\caption{The 25 surveyed studies}\label{tab:survey-overview}\\
\hline \textbf{ID} & \textbf{Study} & \textbf{Venue} & \textbf{Task} \\ \hline \endfirsthead
\hline \textbf{ID} & \textbf{Study} & \textbf{Venue} & \textbf{Task} \\ \hline \endhead
\hline \endfoot
S1 & Hybrid forecasting of commodity prices: machine learning, time series and stochastic simulation~\cite{celik2025} & Borsa \.Istanbul Rev., 2025 & Price, risk \\ \hline
S2 & Multi-commodity price forecasting and anomaly detection with attention-boosted models~\cite{sediyono2025} & J. Agric. Food Res., 2025 & Price, anomaly \\ \hline
S3 & ARIMA--LSTM for volatile prices with random forest~\cite{ray2023} & Appl. Soft Comput., 2023 & Price \\ \hline
S4 & ICEEMDAN--VMD--GRU for economic resilience~\cite{sathe2025} & Results Eng., 2025 & Price \\ \hline
S5 & Volatility dynamics of agricultural futures under uncertainties~\cite{fiszeder2024} & Energy Econ., 2024 & Volatility \\ \hline
S6 & Financial stress and realised volatility of agricultural commodities~\cite{algirdas2024} & Res. Int. Bus. Finance, 2024 & Volatility \\ \hline
S7 & Optimal forecast combination (PSO-CS) for daily futures prices~\cite{zeng2023} & Appl. Soft Comput., 2023 & Price, ensemble \\ \hline
S8 & Transformer anomaly detection with inter-variable attention~\cite{tranad2024} & Inf. Fusion, 2024 & Anomaly \\ \hline
S9 & Hidden-Markov-guided deep models for volatile prices~\cite{avinash2024} & Appl. Soft Comput., 2024 & Price \\ \hline
S10 & Economic policy uncertainty and agricultural prices, China~\cite{epu2024} & Asia Global Econ., 2024 & Volatility \\ \hline
S11 & STL--LSTM for seasonal agricultural price series~\cite{jaiswal2025} & Ann. Data Sci., 2025 & Price \\ \hline
S12 & CEEMDAN hybrid models for agricultural prices~\cite{ceemdan2024} & Sci. Rep., 2024/25 & Price \\ \hline
S13 & Exogenous-variable deep learning for tomato, onion, potato~\cite{exogenous2024} & Sci. Rep., 2024 & Price \\ \hline
S14 & Hybrid LSTM--GARCH volatility forecasting, India~\cite{manogna2025} & J. Big Data, 2025 & Volatility \\ \hline
S15 & Commodity price prediction: a machine-learning framework~\cite{mohanty2023} & Neural Comput. Appl., 2023 & Price \\ \hline
S16 & Demand and price forecasting using deep learning~\cite{demand2023} & LNCS (MIWAI), 2023 & Price, demand \\ \hline
S17 & Point and interval forecasting, decomposition-ensemble and KDE~\cite{kde2024} & Soft Comput., 2024 & Price, interval \\ \hline
S18 & Optimised machine learning for commodity prices~\cite{optimized2023} & Neural Comput. Appl., 2024 & Price \\ \hline
S19 & GCN--BiGRU--TPE futures price prediction~\cite{gcn2024} & Comput. Econ., 2025 & Price, graph \\ \hline
S20 & Crop yield prediction and recommendation, deep ensembles (survey)~\cite{cropyield} & IEEE conf., 2024 & Crop, survey \\ \hline
S21 & Framework for ragi price prediction in Karnataka~\cite{meena2024} & IEEE Access, 2024 & Price \\ \hline
S22 & AgTech: volatility prediction for commodity trading~\cite{agtech2024} & IEEE Access, 2024 & Volatility \\ \hline
S23 & Smart agriculture: state, opportunities and challenges~\cite{smartag2024} & IEEE Access, 2024 & Review \\ \hline
S24 & Smart farming: crop recommendation and price prediction~\cite{smartfarm2025} & IEEE conf., 2025 & Crop, price \\ \hline
S25 & Federated explainable AI for smart agriculture~\cite{fedxai2025} & IEEE Access, 2025 & Crop, XAI \\ \hline
\end{longtable}}

\subsection{Techniques and results}
Tables~\ref{tab:survey-t1} to~\ref{tab:survey-t3} record, for each study, the techniques used and the principal results as reported by the authors.

\Needspace{12\baselineskip}
{\footnotesize\setlength{\tabcolsep}{4pt}\renewcommand{\arraystretch}{1.15}
\begin{longtable}{|p{0.9cm}|p{5.8cm}|p{9.2cm}|}
\caption{Techniques and results, S1--S10 (price, volatility and anomaly studies, Elsevier)}\label{tab:survey-t1}\\
\hline \textbf{ID} & \textbf{Techniques} & \textbf{Principal results} \\ \hline \endfirsthead
\hline \textbf{ID} & \textbf{Techniques} & \textbf{Principal results} \\ \hline \endhead
\hline \endfoot
S1 & ARIMA, VAR, LSTM and stochastic simulation (geometric, fractional and bridge Brownian motion); value-at-risk & LSTM best in high-volatility and crisis periods; ARIMA stable in short-term, low-volatility markets; geometric Brownian motion balances accuracy and risk; fractional motion captures long memory but raises volatility \\ \hline
S2 & Transformer forecasting; attention-boosted LSTM variational autoencoder for anomalies & Transformers beat LSTM and GRU in price forecasting; the attention-boosted autoencoder improves anomaly detection and captures sudden spikes and drops \\ \hline
S3 & ARIMA plus LSTM hybrid; random-forest selection of lags and features; volatility modelling & ARIMA captures linear trend, LSTM the non-linear part, random forest improves feature selection; the hybrid cuts RMSE and MAPE markedly \\ \hline
S4 & ICEEMDAN and VMD decomposition followed by GRU & Decomposition reduces noise and non-stationarity; the hybrid outperforms LSTM and RNN variants \\ \hline
S5 & Heterogeneous autoregressive (HAR) models, LASSO; news-based and market-based uncertainty indicators & Market-based indicators outperform news-based ones; LASSO improves on HAR; uncertainty is strongly linked to price volatility \\ \hline
S6 & HAR-RV variants on high-frequency realised volatility & Financial stress predicts realised volatility in-sample and shows financialisation of agricultural markets; out-of-sample prediction is inconsistent \\ \hline
S7 & EWT, SSA and VMD decomposition; ARIMA, ETS, BPNN, ELM; PSO-CS weight optimisation & The optimised ensemble beats all single models, with a large reduction in MAPE \\ \hline
S8 & Variable temporal transformer with self-attention for multivariate anomaly detection & Captures temporal and inter-variable dependence, beats traditional detectors and gives interpretable anomaly explanations \\ \hline
S9 & Hidden Markov model with RNN, LSTM, GRU, BiLSTM, CNN; technical indicators and FFT & The Markov model captures regime change; guided deep models beat the same models alone for short and medium horizons \\ \hline
S10 & Time-varying-parameter VAR for transmission of policy uncertainty & Policy uncertainty affects price volatility, with effects that vary over time and across commodities \\ \hline
\end{longtable}}

\Needspace{12\baselineskip}
{\footnotesize\setlength{\tabcolsep}{4pt}\renewcommand{\arraystretch}{1.15}
\begin{longtable}{|p{0.9cm}|p{5.8cm}|p{9.2cm}|}
\caption{Techniques and results, S11--S19 (Springer Nature)}\label{tab:survey-t2}\\
\hline \textbf{ID} & \textbf{Techniques} & \textbf{Principal results} \\ \hline \endfirsthead
\hline \textbf{ID} & \textbf{Techniques} & \textbf{Principal results} \\ \hline \endhead
\hline \endfoot
S11 & STL decomposition into trend, seasonal and remainder; an LSTM per component; recombined forecast & Beat ARIMA and a plain LSTM; seasonal components predicted more reliably; RMSE and MAPE improved on benchmarks \\ \hline
S12 & CEEMDAN decomposition with a time-delay neural network per component; compared with ARIMA, random forest, GBM, SVR & Best of all models; RMSE about 20\% lower than other EMD-based models and about 57\% lower than traditional machine learning; MAPE near 2--4\% for major crops \\ \hline
S13 & NBEATSX and TransformerX with weather variables; compared with ARIMAX, ANN, SVR, random forest, XGBoost; RMSE, MAE, sMAPE, MASE & NBEATSX best for all three crops (tomato, onion, potato); deep models clearly beat statistical and machine-learning baselines \\ \hline
S14 & LSTM for long-term pattern; GARCH on LSTM residuals; compared with GARCH, GJR-GARCH, EWMA, MEM & Best for all commodities; onion RMSE fell from 5.44 (GARCH) to 0.70, tomato from 7.67 to 0.79 \\ \hline
S15 & ARIMA, AR, MA, ES, STL for yield; supply from yield, residue and imports; demand by time series; decision tree, RF, KNN, kernel ridge for price & ARIMA(4,1,1) best for yield; decision-tree regressor best for price (RMSE about 0.115, MAPE about 14\%) \\ \hline
S16 & RNN, LSTM and bidirectional LSTM compared with statistical and machine-learning methods; horizon up to one year & LSTM and Bi-LSTM capture temporal dependence better than ARIMA variants; detailed error figures are not available in the accessible preview \\ \hline
S17 & VMD with energy entropy; K-means reconstruction; ARIMA for low and LSTM for high frequencies; WOA-optimised adaptive KDE intervals & Lowest error among compared methods; MAPE about 1--2\% one step ahead; interval coverage (PICP) above 95\% at the 95\% level; 5--15\% lower RMSE than other hybrids \\ \hline
S18 & Extreme learning machine and LSTM tuned by genetic algorithm; Auto-ARIMA; 11 commodities, daily 2000--2022 & GA-ELM best overall (MAPE about 1--3\%), trains faster than GA-LSTM, forecasts 60 days ahead and follows the 2022 peaks \\ \hline
S19 & Multi-graph construction; graph convolution for spatial structure; BiGRU for time; TPE Bayesian tuning & On corn futures, RMSE about 20.2, MAPE about 0.55\%, $R^2$ about 0.97 \\ \hline
\end{longtable}}

\Needspace{12\baselineskip}
{\footnotesize\setlength{\tabcolsep}{4pt}\renewcommand{\arraystretch}{1.15}
\begin{longtable}{|p{0.9cm}|p{5.8cm}|p{9.2cm}|}
\caption{Techniques and results, S20--S25 (IEEE)}\label{tab:survey-t3}\\
\hline \textbf{ID} & \textbf{Techniques} & \textbf{Principal results} \\ \hline \endfirsthead
\hline \textbf{ID} & \textbf{Techniques} & \textbf{Principal results} \\ \hline \endhead
\hline \endfoot
S20 & PSO and WOA for feature selection; SVM, random forest, RNN, Bi-LSTM, CNN, MLP and ensembles (survey of methods) & Ensemble and optimisation approaches outperform traditional methods in crop prediction; Bi-LSTM and CNN strongest on climate and crop data \\ \hline
S21 & LSTM deep-learning framework with temporal and seasonal features; compared with statistical and machine-learning models & Higher accuracy on ragi price trends, robust to non-linear fluctuations, better precision and recall than baselines \\ \hline
S22 & Bi-LSTM on multivariate market indicators (price, volume, open interest); rolling-window volatility index; dropout; binary classification of high/low volatility & Improved accuracy of volatility forecasts for exchange-traded commodities; robust to sudden fluctuations; fewer false predictions \\ \hline
S23 & Survey of IoT, AI and machine learning, blockchain, cloud and precision farming & Documents gains in resource efficiency (water, fertiliser, energy) and applications in precision farming, monitoring and yield prediction \\ \hline
S24 & XGBoost, random forest, decision tree, SVM, KNN, naive Bayes, linear and logistic regression & Integrated crop recommendation and price prediction; ensembles most accurate; price models more reliable than traditional statistics \\ \hline
S25 & Federated learning with explainable AI; CNN, LSTM, RF, SVM, gradient boosting; SHAP and LIME & Better crop-recommendation accuracy across distributed data while preserving privacy; transparent predictions; scales across regions \\ \hline
\end{longtable}}

\subsection{Research gaps and takeaways}
Tables~\ref{tab:survey-g1} to~\ref{tab:survey-g3} record the gaps named for each study in the workbook and the lesson each gives for MandiSense AI.

\Needspace{12\baselineskip}
{\footnotesize\setlength{\tabcolsep}{4pt}\renewcommand{\arraystretch}{1.15}
\begin{longtable}{|p{0.9cm}|p{6.4cm}|p{8.6cm}|}
\caption{Research gaps and takeaways, S1--S10}\label{tab:survey-g1}\\
\hline \textbf{ID} & \textbf{Research gaps} & \textbf{Takeaways for this project} \\ \hline \endfirsthead
\hline \textbf{ID} & \textbf{Research gaps} & \textbf{Takeaways for this project} \\ \hline \endhead
\hline \endfoot
S1 & Models used independently; no shock detection; volatility not in decision support; no explainable architecture & No single model wins in every regime, so a council of specialists is justified; risk-aware, adaptive systems are needed \\ \hline
S2 & Centralised, not agent-based; no explicit risk quantification; limited decision-support integration & Joint forecasting and anomaly detection works; risk and explanation still need to be added \\ \hline
S3 & Monthly data only; no ensemble extension; volatility handled indirectly & Hybrid statistical-plus-deep models beat single models; feature selection matters \\ \hline
S4 & High computational cost; no volatility or anomaly module; one pipeline & Noise reduction is critical for agricultural series; decomposition plus deep learning is effective (and, here, audited for leakage) \\ \hline
S5 & Volatility only, not price; not linked to forecasting or decisions & External uncertainty signals matter; volatility should be tied to price prediction \\ \hline
S6 & Out-of-sample reliability weak; not linked to price models & Financial stress is a volatility driver; regimes change; adaptive models are needed \\ \hline
S7 & Centralised; no explainability or volatility alerts; costly & Combination improves accuracy and stability; weight optimisation is crucial \\ \hline
S8 & Not integrated with price forecasting; little agricultural use; intensive & Attention helps anomaly detection; anomaly detection should be coupled with forecasting \\ \hline
S9 & Commodity- and region-specific; no ensemble or multi-agent design; limited explainability & Regime awareness improves forecasts; probabilistic-plus-deep hybrids suit volatile series \\ \hline
S10 & China only; no forecasting framework & Policy uncertainty drives volatility; time-varying models are needed; policy effects belong in forecasting systems \\ \hline
\end{longtable}}

\Needspace{12\baselineskip}
{\footnotesize\setlength{\tabcolsep}{4pt}\renewcommand{\arraystretch}{1.15}
\begin{longtable}{|p{0.9cm}|p{6.4cm}|p{8.6cm}|}
\caption{Research gaps and takeaways, S11--S19}\label{tab:survey-g2}\\
\hline \textbf{ID} & \textbf{Research gaps} & \textbf{Takeaways for this project} \\ \hline \endfirsthead
\hline \textbf{ID} & \textbf{Research gaps} & \textbf{Takeaways for this project} \\ \hline \endhead
\hline \endfoot
S11 & Weather, supply and demand shocks and market events not used; centralised & Decomposition handles strong seasonality, as in mandi series; preprocessing to isolate components helps deep models \\ \hline
S12 & History only, no weather or demand; centralised & Decomposition before prediction helps volatile prices; hybrids are more reliable than single models \\ \hline
S13 & Only weather; no demand, festivals or exports; centralised & External factors improve accuracy; multi-input forecasting suits Indian markets \\ \hline
S14 & Price history only; centralised & LSTM--GARCH suits abnormal-volatility detection and risk alerts; volatility should join multi-model forecasting \\ \hline
S15 & One crop (corn); no deep learning or real-time data & Integrating supply and demand improves price forecasts; multi-factor beats price alone \\ \hline
S16 & Few metrics reported; no exogenous drivers & Deep sequence models capture temporal trend; test several architectures; report detailed errors \\ \hline
S17 & Price history only; computationally complex & Intervals support risk-aware decisions; ARIMA plus LSTM improves stability; ranges should accompany forecasts \\ \hline
S18 & No external factors; centralised & Optimised hybrids are stable; 60-day forecasts are feasible; compare several models; add arrivals, season and festivals \\ \hline
S19 & One commodity; no weather or demand cycles & Graph structure suits mandis linked by trade; Bayesian tuning automates configuration \\ \hline
\end{longtable}}

\Needspace{12\baselineskip}
{\footnotesize\setlength{\tabcolsep}{4pt}\renewcommand{\arraystretch}{1.15}
\begin{longtable}{|p{0.9cm}|p{6.4cm}|p{8.6cm}|}
\caption{Research gaps and takeaways, S20--S25}\label{tab:survey-g3}\\
\hline \textbf{ID} & \textbf{Research gaps} & \textbf{Takeaways for this project} \\ \hline \endfirsthead
\hline \textbf{ID} & \textbf{Research gaps} & \textbf{Takeaways for this project} \\ \hline \endhead
\hline \endfoot
S20 & Few multi-source datasets; no standard feature engineering; narrow metrics; little interpretability & Ensembles and optimisation help; precision, recall and F1 should be reported; false rates affect trust \\ \hline
S21 & One crop; no socio-economic factors; no deployment study; interpretability & LSTM models suit price series; temporal features help; the framework supports farmers and policymakers \\ \hline
S22 & Policy factors missing; no deployment or multi-commodity scale; interpretability; cross-market links & Volatility prediction can steady trading decisions; ensembles beat single models; external features help \\ \hline
S23 & Little real-time deployment in developing regions; interoperability; privacy; farmer usability & Smart agriculture reduces waste and raises productivity; adoption depends on farmers; solutions must suit smallholders \\ \hline
S24 & No real-time integrated deployment; socio-economic and policy factors; interpretability; scalability & Crop recommendation and price forecasting together support market decisions; multi-source data is crucial \\ \hline
S25 & No real-time rural deployment; interoperability; efficiency against explainability; farmer usability; security & Federated learning avoids sharing raw farmer data; explanations build trust \\ \hline
\end{longtable}}

\subsection{Comparison matrix}
Table~\ref{tab:survey-matrix} compares the studies on ten properties. A tick means the study uses or addresses the property. The last row shows MandiSense AI for reference.

\begin{table}[!t]
\centering
\caption{Comparison of the surveyed studies on ten properties}
\label{tab:survey-matrix}
\footnotesize
\setlength{\tabcolsep}{4.6pt}
\renewcommand{\arraystretch}{1.1}
\begin{tabular}{|l|c|c|c|c|c|c|c|c|c|c|}
\hline
\textbf{ID} & \textbf{Decomp.} & \textbf{RNN} & \textbf{Attn.} & \textbf{Volat.} & \textbf{Ensemble} & \textbf{Exog.} & \textbf{Interval} & \textbf{Anomaly} & \textbf{Graph} & \textbf{Decision} \\ \hline
S1  & & $\checkmark$ & & $\checkmark$ & $\checkmark$ & & & & & \\ \hline
S2  & & & $\checkmark$ & & & & & $\checkmark$ & & \\ \hline
S3  & & $\checkmark$ & & $\checkmark$ & $\checkmark$ & & & & & \\ \hline
S4  & $\checkmark$ & $\checkmark$ & & & & & & & & \\ \hline
S5  & & & & $\checkmark$ & & $\checkmark$ & & & & \\ \hline
S6  & & & & $\checkmark$ & & $\checkmark$ & & & & \\ \hline
S7  & $\checkmark$ & & & & $\checkmark$ & & & & & \\ \hline
S8  & & & $\checkmark$ & & & & & $\checkmark$ & & \\ \hline
S9  & & $\checkmark$ & & $\checkmark$ & & & & & & \\ \hline
S10 & & & & $\checkmark$ & & $\checkmark$ & & & & \\ \hline
S11 & $\checkmark$ & $\checkmark$ & & & & & & & & \\ \hline
S12 & $\checkmark$ & & & & & & & & & \\ \hline
S13 & & & $\checkmark$ & & & $\checkmark$ & & & & \\ \hline
S14 & & $\checkmark$ & & $\checkmark$ & & & & & & \\ \hline
S15 & & & & & & $\checkmark$ & & & & \\ \hline
S16 & & $\checkmark$ & & & & & & & & \\ \hline
S17 & $\checkmark$ & $\checkmark$ & & & $\checkmark$ & & $\checkmark$ & & & \\ \hline
S18 & & $\checkmark$ & & & $\checkmark$ & & & & & \\ \hline
S19 & & $\checkmark$ & & & & & & & $\checkmark$ & \\ \hline
S20 & & $\checkmark$ & & & $\checkmark$ & $\checkmark$ & & & & $\checkmark$ \\ \hline
S21 & & $\checkmark$ & & & & & & & & \\ \hline
S22 & & $\checkmark$ & & $\checkmark$ & & $\checkmark$ & & & & \\ \hline
S23 & & & & & & & & & & \\ \hline
S24 & & & & & $\checkmark$ & $\checkmark$ & & & & $\checkmark$ \\ \hline
S25 & & $\checkmark$ & & & & & & & & $\checkmark$ \\ \hline
\textbf{Total} & \textbf{5} & \textbf{14} & \textbf{3} & \textbf{8} & \textbf{7} & \textbf{8} & \textbf{1} & \textbf{2} & \textbf{1} & \textbf{3} \\ \hline
\textbf{This work} & $\checkmark$ & & $\checkmark$ & $\checkmark$ & $\checkmark$ & & $\checkmark$ & $\checkmark$ & $\checkmark$ & $\checkmark$ \\ \hline
\end{tabular}
\par\smallskip{\footnotesize Decomp.: decomposition; Attn.: attention or transformer; Volat.: volatility or risk modelling; Exog.: external drivers; Interval: prediction intervals; Graph: spatial or graph structure; Decision: output tied to a farmer decision. \emph{This work}: leakage audit of decomposition, zero-shot transformer baseline, volatility desk, council, conformal floors, shock desk, spillover engine, market choice.}
\end{table}

\subsection{Gap analysis across the survey}
The gaps recorded for the 25 studies fall into a small number of themes. Table~\ref{tab:survey-gaps} counts how many studies name each and states how the present work responds.

\begin{table}[!t]
\centering
\caption{Recurring gaps in the surveyed studies and the response of this project}
\label{tab:survey-gaps}
\small
\renewcommand{\arraystretch}{1.15}
\begin{tabular}{|p{4.0cm}|p{3.4cm}|c|p{6.4cm}|}
\hline
\textbf{Gap theme} & \textbf{Studies} & \textbf{n} & \textbf{Response in MandiSense AI} \\ \hline
External drivers (weather, demand, policy, socio-economic) not used & S11--S19, S21, S22, S24 & 12 & Arrivals, season and calendar features in the panel; a decision-focused label; no weather feed (live data not available) \\ \hline
Centralised design, no multi-agent or distributed structure & S2, S3, S7, S9, S11, S12, S14, S18 & 8 & Specialist council with learned weights and per-crop engines behind one API (Section~\ref{sec:architecture}) \\ \hline
Weak interpretability or explainability & S1, S7, S9, S20--S22, S24, S25 & 8 & Why-this-mandi breakdown, what-if, track record and a checked brief (Section~\ref{sec:insight-method}) \\ \hline
No real-time deployment or scalability study & S5, S6, S10, S21--S25 & 8 & A running service, a 60-check smoke test and a morning digest; no field trial yet \\ \hline
Volatility, risk or anomaly not tied to price or decisions & S1, S2, S4, S5, S6, S8, S10 & 7 & Shock desk, volatility view and revenue floors in the same service as the advice \\ \hline
Narrow data (one crop, market or region; monthly only) & S3, S9, S10, S15, S19, S21 & 6 & Four crops, 12 markets, daily data; one region (stated as a limit) \\ \hline
High computational cost & S4, S7, S8, S17 & 4 & Gradient boosting served; heavy models used only as challengers \\ \hline
Judged by forecast error, not by a decision & S1--S4, S7, S9, S11--S19, S21 & 16 & Regret and floor violation as the primary measures (Chapter~\ref{ch:results}) \\ \hline
Decomposition leakage not audited (our reading) & S4, S7, S11, S12, S17 & 5 & Walk-forward re-decomposition audit (Section~\ref{sec:leakage}) \\ \hline
\end{tabular}
\end{table}

Three conclusions follow. First, hybrid and decomposition-based models dominate recent work (14 of 25 use recurrent networks, 5 use decomposition) and are judged almost entirely by forecast error. Second, only one study (S17) provides an interval and none states a guarantee or a downside in rupees. Third, no study treats the farmer's own choice, of market and day, as the quantity to be optimised, and no study audits how much of the reported gain comes from fitting a decomposition on the whole series. These are the gaps that Section~\ref{sec:limitations} turns into objectives.
